\uuid{4E7y}
\titre{Primitive}
\auteur{bodin ; croizat ; sacre}
\organisation{exo7}
\licence{CC BY-NC-SA 4.0}
\datecreate{2022-03-08}
\niveau{L1}
\module{Analyse}
\chapitre{Calcul d'intégrales}
\sousChapitre{Primitives diverses}
\difficulte{2}
\theme{Primitive}
\exo7qcmid{hilisit.equations-differentielles.005}

\contenu{
\begin{question}
\qclassification{127.11}
On considère la fonction \(f:x\mapsto 2 e^{-2x}-3\). Quelles sont les affirmations exactes ?
\begin{answers}
  \bad{\(f\) est une primitive de \(- e^{-2x}-3x\) sur \(\Rr\).}
  \good{\(f\) est une primitive de \(-4 e^{-2x}\) sur \(\Rr\).}
  \good{\(f\) est la primitive de \(-4 e^{-2x}\) sur \(\Rr\) valant \(-1\) en \(x=0\).}
  \bad{\(f\) est la dérivée de \(x\mapsto - e^{-2x}\)}
\end{answers}
\begin{explanations}
Pour vérifier si une fonction \(f\) est une primitive d'une fonction \(g\), on calcule la dérivée de  \(f\) et on regarde si on obtient bien la fonction \(g\). La dérivée de $ e^{-2x}$ est $-2 e^{-2x}$ donc $f'(x)=-4 e^{-2x}$. De plus $f(0)=2 e^0-3=2-3=-1$.
\end{explanations}
\end{question}
}
